\section{Discussão}\label{sec:discussao}

\textbf{On-off.} Sem um nível igual ao esterçamento de regime de cada curva, o on-off nunca atinge o equilíbrio e o erro descreve um ciclo-limite. Sua amplitude depende do tempo de reversão do comando (histerese, período de $50\,\mathrm{ms}$ e limite de taxa, $\approx0{,}5\,\mathrm{s}$ de $+15^\circ$ a $-15^\circ$), durante o qual o veículo se desloca lateralmente a uma taxa proporcional a $v$. Por isso o erro cresce com a velocidade e diverge a $2{,}5\,\mathrm{m/s}$ ainda na reta. Com a câmera, que enxerga apenas $\pm30^\circ$ e não vê os primeiros $0{,}5\,\mathrm{m}$ à frente do eixo, uma oscilação que coincide com uma curva fechada leva a linha para fora da imagem. A amplitude $\delta_{on}$ impõe um compromisso entre o tamanho do ciclo-limite e o raio mínimo de curva, $L/\tan\delta_{on}$. A distância de antecipação $L_d$ foi essencial: sem o termo $\arctan(e_y/L_d)$, o on-off corrigiria apenas a orientação.

\textbf{Stanley.} Em regime, o termo $e_\psi$ fornece o esterçamento necessário para a curvatura, e o erro fica na faixa de milímetros com referências ideais e de $1$--$4\,\mathrm{cm}$ com a câmera. Os maiores erros com visão aparecem nos cantos vivos da linha pintada: o ajuste local, feito com pontos a mais de $0{,}5\,\mathrm{m}$ à frente, é extrapolado até o eixo e não representa a quina, e o atuador limitado em taxa não reproduz o salto de esterçamento. Por isso o pico cresce com a velocidade (de $6$ para $33\,\mathrm{cm}$) enquanto o RMS cresce pouco, e a redução de velocidade nas curvas diminui o pico pela metade com pequeno custo no tempo de volta.

\textbf{Do simulador para o veículo real.} A mesma lei, com os mesmos ganhos, manteve o carrinho real sobre a faixa com erro estimado de $2$--$3\,\mathrm{cm}$ a $\approx0{,}6\,\mathrm{m/s}$. Os erros médios positivos, com esterçamento médio de $6$--$7^\circ$ numa faixa praticamente reta, indicam um viés sistemático, como desalinhamento mecânico da direção ou erro de guinada na calibração da câmera. Como o Stanley não possui ação integral, um desvio constante só é compensado com erro permanente; uma correção de \textit{trim} ou um termo integral pequeno eliminariam esse viés. O veículo real ainda apresenta atraso visual de $\approx0{,}15\,\mathrm{s}$, perdas de detecção e falhas de outros subsistemas (encoder, ultrassom) que interromperam ensaios. O on-off não foi testado no veículo real; pela simulação, ele exigiria velocidades baixas.

\textbf{Limitações.} No simulador, a projeção e a memória usam a pose da câmera fornecida pela cena, o que equivale a odometria ideal, e a segmentação HSV foi avaliada com iluminação constante. O erro verdadeiro tem resolução de $1{,}8\,\mathrm{cm}$, de modo que valores abaixo de $\approx1\,\mathrm{cm}$ estão próximos do limite da referência, e cada condição foi ensaiada uma única vez. No veículo real não há referência independente: os erros são estimados pela própria câmera, a odometria (encoder e modelo cinemático, sem IMU) acumula deriva e não foi usada como medida, e apenas trechos de $30\,\mathrm{s}$ foram percorridos.
